\chapter{Variational Quantum Algorithms and QAOA}
\label{ch:lec18}

\begin{objectives}
\begin{itemize}[leftmargin=1.4em,nosep]
  \item State the variational principle behind VQE.
  \item Construct the QAOA ansatz for combinatorial optimization.
  \item Relate QAOA to adiabatic evolution.
  \item Discuss trainability and barren-plateau concerns.
\end{itemize}
\end{objectives}

\todo{Write the introductory paragraph orienting the lecture.}

\section{The variational quantum eigensolver (VQE)}
\label{sec:lec18-1}
\todo{Draft this section.}

\section{The quantum approximate optimization algorithm (QAOA)}
\label{sec:lec18-2}
\todo{Draft this section.}

\section{The adiabatic connection}
\label{sec:lec18-3}
\todo{Draft this section.}

\section{Trainability and expressibility}
\label{sec:lec18-4}
\todo{Draft this section.}

\exercisesection
\begin{exercise}
\todo{Add exercise.}
\end{exercise}
\begin{exercise}
\todo{Add exercise.}
\end{exercise}
\begin{exercise}
\todo{Add exercise.}
\end{exercise}

\readingsection
\begin{itemize}[leftmargin=1.4em]
  \item E.~Farhi, J.~Goldstone, S.~Gutmann, \emph{A Quantum Approximate Optimization Algorithm}, arXiv:1411.4028.
  \item M.~Cerezo et al., \emph{Variational Quantum Algorithms}, Nat.~Rev.~Phys.~3 (2021).
\end{itemize}
